%=============================================================================
% Wave 3, Iteration 4: Patent 15 (Generator / SRF / Superconducting Materials)
%
% FOCUS: Alternative superconducting materials beyond Nb3Sn+YBCO+La3Ni2O7
%        MgB2, infinite-layer nickelate (NdNiO2), heavy fermions (CeCoIn5,
%        UTe2), flat-band materials (MATBG, twisted WSe2/MoS2).
%        Full FOM derivation, comparison table, updated resonator recommendation.
%
% Agent: P15 Generator Cross-Reference (Iteration 4)
% Date:  March 2026
%=============================================================================

\documentclass[11pt]{article}
\usepackage{amsmath,amssymb,amsthm}
\usepackage{geometry}
\geometry{margin=1in}
\usepackage{booktabs}
\usepackage{enumitem}
\usepackage{hyperref}
\usepackage{xcolor}
\usepackage{siunitx}
\usepackage{float}
\usepackage{mdframed}
\usepackage{multirow}
\usepackage{array}
\usepackage{bm}

\newtheorem{theorem}{Theorem}
\newtheorem{proposition}{Proposition}
\newtheorem{lemma}{Lemma}
\newtheorem{finding}{Finding}
\newtheorem{correction}{Correction}

\newcommand{\ee}[1]{\times 10^{#1}}
\newcommand{\psifield}{\psi}
\newcommand{\Opsi}{\Omega_\psi}
\newcommand{\gpsi}{\gamma_\psi}
\newcommand{\Mpsi}{M_\psi}
\newcommand{\lam}{\lambda}
\newcommand{\order}[1]{\mathcal{O}(#1)}
\newcommand{\boxedfinding}[1]{\begin{mdframed}[backgroundcolor=yellow!10]#1\end{mdframed}}
\newcommand{\lambdaL}{\lambda_{\rm L}}
\newcommand{\xsc}{\xi_{\rm sc}}
\newcommand{\Tc}{T_{\rm c}}
\newcommand{\kB}{k_{\rm B}}
\newcommand{\FOM}{\mathrm{FOM}}
\newcommand{\mstar}{m^*}

\title{Wave 3 Iteration 4: Patent 15 Generator ---\\
Alternative Superconducting Materials for DFD Coupling\\
\large MgB$_2$, Infinite-Layer Nickelate (NdNiO$_2$), Heavy Fermions,\\
Flat-Band Materials: Full FOM Derivation and Ranking}
\author{DFD Research Programme --- P15 Cross-Reference Agent (Iteration 4)}
\date{March 2026}

\begin{document}
\maketitle
\tableofcontents
\newpage

%=============================================================================
\section{W3i4 P15: Iteration Context and Inherited Corrections}
\label{sec:context}
%=============================================================================

This document is Wave~3, Iteration~4 of the P15 analysis. It carries forward
all corrections from W3i1--W3i3 and performs the mandated extension to
alternative superconducting materials not examined in prior iterations.

\subsection{Consolidated inherited corrections (W3i1--W3i3)}

\begin{enumerate}
\item \textbf{SQMS bound (W3i2)}: The naive $|\lam-1| < 7\ee{-10}$ SQMS
  constraint is \emph{invalidated} by BCS gap suppression of quasiparticle
  conductivity (Process B suppressed by $e^{-\Delta/\kB T} \sim 10^{-28}$).
  The SQMS $Q$-data provides no valid upper limit on $|\lam-1|$.

\item \textbf{YBCO enhancement (W3i2)}: Bulk anisotropic coupling acts only
  within the London penetration depth, giving $< 10^{-3}\times$ bulk
  enhancement. The surface mechanism (differential $\lambdaL$: YBCO
  $150$~nm vs.\ Nb $40$~nm) gives $4.0\times$ surface enhancement.

\item \textbf{Process A / B distinction (W3i3)}: P15 operates via Process A
  (EM $\to \psi$ sourcing through $\square\psi = -4\pi G T/c^4$). This is
  independent of quasiparticle occupation and unaffected by the BCS gap.
  Process B (thermalization through quasiparticle coupling) is suppressed
  by $\sim 10^{-28}$.

\item \textbf{Virtual Cooper pair fluctuations (W3i3)}: Contribution
  $\psi_{\rm fluct} \sim 10^{-48}$; negligible.

\item \textbf{Andreev channel (W3i3)}: Suppressed by same exponential factor
  as quasiparticle thermal occupation; negligible.

\item \textbf{Optimal phase gate (W3i3)}: $|\alpha/\beta|_\star = 1 + \sqrt{2}$;
  balanced superposition gives zero signal.

\item \textbf{La$_3$Ni$_2$O$_7$ (W3i3 MatSci)}: Mixed $d_{x^2-y^2} + s_\pm$
  pairing confirmed; $\xsc = 3.15$~nm; Figure of Merit
  $\FOM_{\rm RT} = 252$~K (defined in W3i3 MatSci, see Section~\ref{sec:fom}
  for revised definition).

\item \textbf{Nb$_3$Sn + YBCO laminate (W3i3)}: Identified as optimal
  resonator geometry. This recommendation is re-evaluated in
  Section~\ref{sec:recommendation}.

\item \textbf{MATBG flat-band enhancement (W3i3)}: $\sim 2700\times$ DoS
  amplification relative to conventional metal.

\item \textbf{Topological superconductors (W3i3)}: Majorana zero modes (MZMs)
  are NOT lifted by DFD coupling; degeneracy is protected by topology.
\end{enumerate}

%=============================================================================
\section{DFD Process A Coupling: Figure of Merit Derivation}
\label{sec:fom}
%=============================================================================

\subsection{Source equation and coupling to material}

The DFD wave equation for the $\psi$-field sourced by an electromagnetic
cavity is:
\begin{equation}
\square\psi = -\frac{4\pi G}{c^4}\,T^{\mu}_{\ \mu},
\label{eq:wave}
\end{equation}
where $T^\mu_{\ \mu} = \rho c^2 - 3P$ is the stress-energy trace. For an
electromagnetic field in a resonant cavity:
\begin{equation}
T^{\rm EM}_{\mu}^{\ \mu} = 0
\quad \text{(traceless for a pure EM field)}.
\label{eq:em-trace}
\end{equation}
The $\psi$-source therefore arises from the \emph{material} response to the
EM field. In the cavity walls, the conduction electrons respond to surface
currents with a kinetic energy density:
\begin{equation}
u_{\rm kin} = \frac{n_s m v_s^2}{2}
= \frac{J_s^2}{2 n_s e^2 / m},
\end{equation}
where $n_s$ is the superfluid density and $v_s = J_s/(n_s e)$ is the
superfluid velocity. The stress-energy contribution from this kinetic
energy over the London penetration depth $\lambdaL$ is:
\begin{equation}
T_{\rm wall}^{\mu}_{\ \mu} = u_{\rm kin} \cdot \lambdaL
\quad [\text{per unit surface area}].
\end{equation}

\subsection{Surface-current-driven $\psi$-source}

The surface current density in an SRF cavity is $J_s = H_{\rm peak} / \lambdaL$,
giving:
\begin{equation}
u_{\rm kin} = \frac{J_s^2 \lambdaL^2 m}{2 n_s e^2 \lambdaL^2}
= \frac{\mu_0 H_{\rm peak}^2}{2}
\cdot \frac{m}{n_s e^2 \mu_0 \lambdaL^2}.
\end{equation}
Using the London relation $\lambdaL^2 = m / (\mu_0 n_s e^2)$:
\begin{equation}
u_{\rm kin} = \frac{\mu_0 H_{\rm peak}^2}{2}
= \frac{B_{\rm peak}^2}{2\mu_0},
\end{equation}
which is simply the stored EM energy density. The surface-integrated
stress-energy source driving the $\psi$-field is then:
\begin{equation}
S_\psi \equiv \int_{\rm walls} T_{\rm wall}\,dA
\propto B_{\rm peak}^2 \cdot A_{\rm cav} \cdot \lambdaL.
\label{eq:source}
\end{equation}

The key result from Eq.~(\ref{eq:source}) is that the $\psi$-source scales
as $\lambdaL$: \emph{materials with larger London penetration depths source
more $\psi$-field per unit surface area at fixed peak field $B_{\rm peak}$.}

\subsection{Process A FOM derivation}

The $\psi$-field amplitude at the cavity centre from wall sourcing is:
\begin{equation}
\psi_{\rm cavity} \propto \frac{G}{c^4}\,S_\psi\,\frac{1}{\omega^2 / c^2}
\propto \lambdaL \cdot \frac{1}{\omega_{\rm cav}^2}.
\end{equation}

For fixed cavity frequency and peak field, the relevant material-dependent
factor is simply $\lambdaL$. However, the material must:
\begin{enumerate}
\item Support high $B_{\rm peak}$ before flux entry (limited by $B_{c1}$ for
  type~II, $B_c$ for type~I);
\item Maintain low RF surface resistance $R_s$ so that the cavity $Q$ is not
  degraded (which would reduce stored energy);
\item Be practically fabricable as a cavity wall material.
\end{enumerate}

The surface resistance for a BCS superconductor at temperature $T \ll \Tc$ is:
\begin{equation}
R_s^{\rm BCS} \approx A\,\frac{\omega^2}{\Tc}\,e^{-1.76\,\Tc/T},
\label{eq:Rs}
\end{equation}
where $A$ is a material-dependent prefactor. Higher $\Tc$ gives exponentially
lower $R_s$ at fixed operating temperature.

The maximum achievable stored energy density before quench is
$u_{\rm max} \propto B_{c2}^2$ (for type~II), while the critical field
scales approximately as $B_{c2} \approx \Phi_0 / (2\pi \xsc^2)$ where
$\Phi_0$ is the flux quantum and $\xsc$ is the coherence length.

Combining all factors, the DFD Process~A Figure of Merit for SRF cavity
materials is:

\begin{equation}
\boxed{
\FOM_{\rm DFD} \equiv \frac{\lambdaL}{\xsc} \cdot \Tc^{1/2}
= \kappa \cdot \Tc^{1/2}
}
\label{eq:fom-dfd}
\end{equation}

where $\kappa = \lambdaL / \xsc$ is the Ginzburg-Landau parameter. The
rationale for each factor:

\begin{itemize}
\item $\lambdaL$: directly enhances $\psi$-source from surface kinetic
  energy (Eq.~\ref{eq:source});
\item $\xsc^{-1}$: $B_{c2} \propto \xsc^{-2}$, so larger $\kappa$ implies
  higher field tolerance; also, smaller $\xsc$ gives higher Cooper pair
  density fluctuations (relevant for quantum corrections);
\item $\Tc^{1/2}$: exponential BCS suppression of surface resistance
  means higher $\Tc$ gives lower losses and higher stored $B$-field at
  fixed operating $T$; the $1/2$ exponent comes from the prefactor of
  Eq.~(\ref{eq:Rs}) at fixed $T/\Tc$.
\end{itemize}

\begin{mdframed}[backgroundcolor=gray!5]
\textbf{Note on alternative FOM form.} The W3i3 MatSci agent used
$\FOM_{\rm RT} \equiv \Tc \cdot \lambdaL^{-2} \cdot \xsc^{-1}$ as a
``room-temperature usability'' proxy. That FOM weights superfluid density
(via $\lambdaL^{-2}$) rather than penetration depth, and is appropriate when
evaluating materials for \emph{shielding} applications. For DFD Process~A
cavity sourcing, where we want \emph{maximum surface kinetic energy
contribution}, the correct FOM is $\kappa\,\Tc^{1/2}$ (Eq.~\ref{eq:fom-dfd}).
Both FOMs are tabulated in Table~\ref{tab:fom} for comparison.
\end{mdframed}

%=============================================================================
\section{MgB$_2$ Evaluation}
\label{sec:mgb2}
%=============================================================================

\subsection{Material parameters}

Magnesium diboride (MgB$_2$) is a two-band phonon-mediated superconductor
with $\Tc = 39$~K, discovered in 2001. Its key parameters:

\begin{itemize}
\item $\Tc = 39$~K (ambient pressure, bulk);
\item Two-band structure: $\sigma$-band ($p_{x,y}$ boron orbitals, strong
  in-plane $E_{2g}$ phonon coupling, $\Delta_\sigma \approx 7.1$~meV) and
  $\pi$-band ($p_z$ boron orbitals, weaker coupling, $\Delta_\pi \approx 2.3$~meV);
\item London penetration depth: $\lambdaL^{ab} \approx 85$~nm ($\sigma$-band
  dominated, in-plane);
\item Coherence length: $\xsc^{ab} \approx 3.5$--$5$~nm (in-plane),
  $\xsc^{c} \approx 15$--$20$~nm (out-of-plane);
\item Ginzburg-Landau parameter: $\kappa_{ab} = \lambdaL^{ab}/\xsc^{ab}
  \approx 85/4.25 \approx 20$;
\item $B_{c2}^{ab} \approx 16$--$18$~T (upper critical field, in-plane);
\item $B_{c1} \approx 30$~mT.
\end{itemize}

\subsection{DFD FOM calculation for MgB$_2$}

Using the definition of Eq.~(\ref{eq:fom-dfd}) with $\lambdaL = 85$~nm,
$\xsc = 4.25$~nm (geometric mean of reported range $3.5$--$5$~nm),
$\Tc = 39$~K:
\begin{equation}
\kappa_{\rm MgB_2} = \frac{85}{4.25} \approx 20.0,
\end{equation}
\begin{equation}
\FOM_{\rm DFD}^{\rm MgB_2}
= 20.0 \times \sqrt{39}
= 20.0 \times 6.24
\approx 125~\text{(dimensionless, in units of nm $\cdot$ K$^{1/2}$/nm)}.
\label{eq:fom-mgb2}
\end{equation}

\subsection{Two-band effects on Process A}

The two-band nature of MgB$_2$ introduces a subtlety for DFD Process~A
coupling. The stress-energy trace from surface kinetic energy involves the
total superfluid density $n_s = n_s^\sigma + n_s^\pi$. Since the London depth
is determined by $\lambdaL^{-2} \propto n_s$, the \emph{total} superfluid
density contributes to the electromagnetic response but the $\sigma$-band
(which has the longer coherence length in the $c$-direction and the larger gap)
dominates the in-plane $\lambdaL$.

\begin{finding}[MgB$_2$ Two-Band DFD Coupling]
\label{finding:mgb2-two-band}
The two-band nature does not create a new DFD coupling channel. Both
$\sigma$ and $\pi$ bands contribute to Process~A through their respective
superfluid densities. The effective $\lambdaL^{\rm eff}$ for DFD purposes
is the measured in-plane value of $\approx 85$~nm. There is no anomalous
enhancement from inter-band scattering: the two-gap structure is irrelevant
to the EM stress-energy trace.
\end{finding}

\subsection{Process A coupling: comparison to Nb, Nb$_3$Sn, YBCO}

At fixed peak surface field $B_{\rm peak} = 200$~mT (a representative SRF
operating point), the surface kinetic energy density is:
\begin{align}
u_{\rm kin}^{\rm Nb} &\propto \lambdaL^{\rm Nb} = 40~\text{nm}, \\
u_{\rm kin}^{\rm Nb_3Sn} &\propto \lambdaL^{\rm Nb_3Sn} = 65~\text{nm}, \\
u_{\rm kin}^{\rm MgB_2} &\propto \lambdaL^{\rm MgB_2} = 85~\text{nm}, \\
u_{\rm kin}^{\rm YBCO} &\propto \lambdaL^{\rm YBCO} = 150~\text{nm}.
\end{align}

The Process~A sourcing advantage of MgB$_2$ over Nb is $85/40 = 2.1\times$,
while YBCO still leads at $150/40 = 3.75\times$. MgB$_2$ lies between
Nb$_3$Sn and YBCO.

\subsection{Practical fabricability of MgB$_2$}

\begin{itemize}
\item \textbf{Thin film deposition}: MgB$_2$ films are grown by thermal
  co-evaporation, HPCVD (hybrid physical-chemical vapour deposition), or
  PLD. Film quality is sensitive to stoichiometry and oxygen contamination.
  Critical temperatures of thin films approach bulk values only for high-purity
  films grown at $\sim 650$~$^\circ$C. This temperature is compatible with
  Nb substrate processing.
\item \textbf{RF surface resistance}: MgB$_2$ thin films achieve
  $R_s \approx 1$--$5~\text{n}\Omega$ at 1.5~GHz and 4~K. This is comparable
  to Nb$_3$Sn at the same temperature ratio $T/\Tc$.
\item \textbf{Operating temperature}: $\Tc = 39$~K allows operation at
  $\sim 15$--$20$~K (closed-cycle cryo), eliminating liquid helium. This
  is a major practical advantage over Nb ($\Tc = 9.2$~K, must operate at
  $\sim 2$~K) and Nb$_3$Sn ($\Tc = 18$~K, typically $\sim 4$~K).
\item \textbf{Coating existing cavities}: MgB$_2$ can in principle be
  deposited as a thin film on Nb or Cu cavity substrates, similar to
  Nb$_3$Sn. Layer thickness $\sim 1$--$3~\mu$m is sufficient to exceed
  $\lambdaL$ and $B_{c1}$ performance.
\end{itemize}

\begin{finding}[MgB$_2$ DFD Assessment]
\label{finding:mgb2}
MgB$_2$ achieves $\FOM_{\rm DFD} \approx 125$, placing it between
Nb$_3$Sn ($\FOM \approx 94$) and YBCO ($\FOM \approx 285$). Its practical
advantage is 15--20~K operation without liquid helium. For a P15 generator
targeting sustained CW operation, MgB$_2$ on a Nb substrate offers a
technically simpler thermal management path than the Nb$_3$Sn+YBCO laminate,
at the cost of $\sim 2\times$ reduced Process~A sourcing relative to YBCO.
\end{finding}

%=============================================================================
\section{Infinite-Layer Nickelate (NdNiO$_2$) Evaluation}
\label{sec:ndnio2}
%=============================================================================

\subsection{Material parameters}

NdNiO$_2$ is the infinite-layer nickelate (ILN) with $d^9$ Ni valence,
structurally analogous to CaCuO$_2$ (the parent compound of cuprates). Key
parameters:

\begin{itemize}
\item $\Tc = 9$--$15$~K (thin film on SrTiO$_3$ substrate, $\sim 20$\%
  Sr doping): $\rm Nd_{0.8}Sr_{0.2}NiO_2$;
\item Under pressure $\sim 14$~GPa: $\Tc$ may reach $80$~K (some reports),
  but reproducibility is disputed. Ambient-pressure value $\Tc \approx 12$~K
  is adopted here as established;
\item Crystal structure: square NiO$_2$ planes, no apical oxygen (unlike
  cuprates and La$_3$Ni$_2$O$_7$);
\item Electronic structure: Ni $3d^9$ + rare-earth $5d$ self-doping band.
  The Fermi surface has $\alpha$ (Ni $d_{x^2-y^2}$) and $\gamma$ (Nd $5d$)
  sheets;
\item London penetration depth: $\lambdaL \approx 200$--$270$~nm (from
  muSR and THz optical measurements on thin films); the large value reflects
  low superfluid density due to dilute doping;
\item Coherence length: $\xsc \approx 3$--$5$~nm (from $B_{c2}$ measurements);
\item Pairing symmetry: controversial --- $d_{x^2-y^2}$ strongly suggested
  by theoretical calculations (DFT+DMFT, fRG), but experimental evidence
  is limited. The DFD framework (W3i3) predicts $d$-wave for single-layer
  infinite-layer structure (no interlayer channel to drive $s_\pm$).
\end{itemize}

\subsection{Comparison to La$_3$Ni$_2$O$_7$}

\begin{center}
\begin{tabular}{lcc}
\toprule
Parameter & NdNiO$_2$ (ILN) & La$_3$Ni$_2$O$_7$ (bilayer) \\
\midrule
$\Tc$ (K) & $12$ & $80$ (at 30 GPa) \\
$\lambdaL$ (nm) & $230$ (average) & $\sim 100$ (estimated) \\
$\xsc$ (nm) & $4$ & $3.15$ \\
$\kappa$ & $57.5$ & $\sim 32$ \\
Pairing & $d_{x^2-y^2}$ (DFD pred.) & $d + s_\pm$ mixed \\
\bottomrule
\end{tabular}
\end{center}

NdNiO$_2$ has a \emph{larger} $\lambdaL$ than La$_3$Ni$_2$O$_7$, which
increases Process~A sourcing, but its $\Tc$ is much lower. The FOM:
\begin{equation}
\kappa_{\rm NdNiO_2} = \frac{230}{4} \approx 57.5,
\end{equation}
\begin{equation}
\FOM_{\rm DFD}^{\rm NdNiO_2}
= 57.5 \times \sqrt{12}
= 57.5 \times 3.46
\approx 199.
\label{eq:fom-ndnio2}
\end{equation}

\subsection{Superfluid density comparison}

The superfluid density $\rho_s \propto \lambdaL^{-2}$:
\begin{align}
\rho_s^{\rm NdNiO_2} &\propto (230~\text{nm})^{-2} = 1.9\ee{-5}~\text{nm}^{-2}, \\
\rho_s^{\rm La327} &\propto (100~\text{nm})^{-2} = 1.0\ee{-4}~\text{nm}^{-2}.
\end{align}

La$_3$Ni$_2$O$_7$ has $\sim 5\times$ higher superfluid density than NdNiO$_2$.
This means YBCO or La$_3$Ni$_2$O$_7$ provide better shielding and better
Meissner-effect field exclusion, while NdNiO$_2$'s larger $\lambdaL$ makes
it a better Process~A source.

\subsection{Practical assessment: thick film growth}

Can infinite-layer nickelates be grown as thick films ($>1$~mm) needed for
cavity walls? The answer is \textbf{no}, at present:

\begin{itemize}
\item All ILN superconductivity has been demonstrated only in epitaxial thin
  films ($\sim 5$--$20$~nm) on SrTiO$_3$ or LaAlO$_3$ substrates. The
  SrTiO$_3$ substrate provides compressive strain that is apparently essential
  for superconductivity;
\item The infinite-layer structure is thermodynamically metastable: bulk
  NdNiO$_2$ has not been synthesised. The bulk parent compound NdNiO$_3$
  must be topotactically reduced using CaH$_2$, and this process only works
  for thin films;
\item Thick films ($> 100$~nm) have not been demonstrated with reproducible
  $\Tc$ values. The strain relaxes above $\sim 15$--$20$~nm, suppressing
  superconductivity.
\end{itemize}

\begin{finding}[NdNiO$_2$ Practical Barrier]
\label{finding:ndnio2-practical}
Infinite-layer nickelate NdNiO$_2$ cannot be fabricated as cavity-wall
material with current technology. Its superconductivity requires epitaxial
thin films ($< 20$~nm) on specific substrates. This is a fundamental barrier,
not merely an engineering challenge. NdNiO$_2$ is excluded from practical
P15 resonator consideration unless a bulk synthesis route is discovered.
Its high FOM$_{\rm DFD} \approx 199$ is thus inaccessible in the near term.
\end{finding}

%=============================================================================
\section{Heavy Fermion Superconductors}
\label{sec:hf}
%=============================================================================

\subsection{Background: effective mass and DFD coupling}

Heavy fermion (HF) superconductors have quasiparticle effective masses
$\mstar \gg m_e$, arising from hybridisation of $f$-electrons with conduction
bands (Kondo lattice physics). For DFD Process~A, the relevant quantity is
the electromagnetic response of the material, which enters through the
superfluid density:
\begin{equation}
n_s = \frac{n e^2}{\mstar c^2} \cdot \frac{\lambdaL^{-2} c^2}{4\pi}
\quad \Rightarrow \quad
\lambdaL^2 = \frac{\mstar c^2}{4\pi n e^2}.
\end{equation}

Large $\mstar$ means large $\lambdaL$, which enhances Process~A
(Eq.~\ref{eq:source}). However, the critical temperature $\Tc$ is typically
very low ($\ll 10$~K) for heavy fermions. The DFD FOM (Eq.~\ref{eq:fom-dfd})
penalises low $\Tc$ through the $\Tc^{1/2}$ factor.

The key question is whether the anomalously large $\lambdaL$ of heavy fermions
compensates for their very low $\Tc$.

\subsection{CeCoIn$_5$}

CeCoIn$_5$ is a heavy fermion superconductor with:
\begin{itemize}
\item $\Tc = 2.3$~K;
\item Sommerfeld coefficient $\gamma \approx 300$~mJ/(mol$\cdot$K$^2$)
  (``heavy'' by conventional SC standards: $\gamma_{\rm Nb} \approx 7$~mJ/mol$\cdot$K$^2$);
\item Effective mass $\mstar / m_e \approx 50$--$100$ (estimated from $\gamma$);
\item London penetration depth: $\lambdaL \approx 200$--$350$~nm (anisotropic;
  in-plane $\approx 240$~nm);
\item Coherence length: $\xsc \approx 5$--$8$~nm;
\item Pairing: $d_{x^2-y^2}$ (nodal, confirmed by specific heat and NMR);
\item $\kappa \approx 240/6 = 40$.
\end{itemize}

DFD FOM:
\begin{equation}
\FOM_{\rm DFD}^{\rm CeCoIn_5}
= 40 \times \sqrt{2.3}
= 40 \times 1.52
\approx 61.
\label{eq:fom-cecoin5}
\end{equation}

The low $\Tc$ severely penalises CeCoIn$_5$ despite its large $\kappa$.

\subsection{Does large $m^*$ enhance Process A?}

The Process~A coupling coefficient scales as:
\begin{equation}
C_\psi^{\rm Process A} \propto \lambdaL \cdot B_{\rm peak}^2,
\end{equation}
and $\lambdaL \propto (\mstar)^{1/2}$. Therefore large $\mstar$ does enhance
Process~A coupling \emph{per unit area at fixed peak field}. However, the
large effective mass also reduces $B_{c2}$:
\begin{equation}
B_{c2} = \frac{\Phi_0}{2\pi \xsc^2}, \quad \xsc \propto v_F / \Tc.
\end{equation}
Since $v_F \propto (\mstar)^{-1}$, we have $\xsc \propto (\mstar \Tc)^{-1}$,
giving $B_{c2} \propto \mstar^2 \Tc^2$. This is favourable, but the
operating temperature must be $T \ll \Tc$, and with $\Tc = 2.3$~K this
requires millikelvin cooling --- entirely impractical for P15.

\begin{finding}[Heavy Fermion CeCoIn$_5$: Impractical for P15]
\label{finding:cecoin5}
CeCoIn$_5$ achieves $\FOM_{\rm DFD} \approx 61$, substantially below
Nb$_3$Sn, YBCO, and MgB$_2$. More critically, its operating temperature
($T < 1$~K for low $R_s$) requires dilution refrigeration, making it
entirely unsuitable for a P15-class SRF generator. The large effective mass
enhances $\lambdaL$ (beneficial for Process~A) but this is overwhelmed by the
extreme penalty on $\Tc$.
\end{finding}

\subsection{UBe$_{13}$}

UBe$_{13}$ parameters:
\begin{itemize}
\item $\Tc = 0.86$~K;
\item $\gamma \approx 1000$~mJ/(mol$\cdot$K$^2$) (among the highest known);
\item $\mstar/m_e \approx 200$--$300$;
\item $\lambdaL \approx 800$~nm (estimated from large $\mstar$);
\item $\xsc \approx 10$--$20$~nm (estimated);
\item Pairing symmetry: unclear (possibly spin-triplet, possible two
  superconducting phases).
\end{itemize}

DFD FOM:
\begin{equation}
\kappa_{\rm UBe_{13}} \approx \frac{800}{15} \approx 53,
\end{equation}
\begin{equation}
\FOM_{\rm DFD}^{\rm UBe_{13}}
= 53 \times \sqrt{0.86}
= 53 \times 0.93
\approx 49.
\label{eq:fom-ube13}
\end{equation}

Even lower than CeCoIn$_5$ due to the very low $\Tc$. Not viable for P15.

\subsection{UTe$_2$ and topological superconductivity}

UTe$_2$ is a spin-triplet superconductor with:
\begin{itemize}
\item $\Tc = 1.6$~K;
\item Strong evidence for spin-triplet $p$-wave pairing (NMR, polar Kerr
  effect, edge-state conductance);
\item Topological superconductor candidate: Chern number $> 0$ on certain
  surfaces;
\item Estimated $\lambdaL \approx 500$~nm, $\xsc \approx 10$--$15$~nm;
\item $\kappa \approx 500/12 \approx 42$;
\item Sommerfeld coefficient $\gamma \approx 120$~mJ/(mol$\cdot$K$^2$).
\end{itemize}

DFD FOM:
\begin{equation}
\FOM_{\rm DFD}^{\rm UTe_2}
= 42 \times \sqrt{1.6}
= 42 \times 1.26
\approx 53.
\label{eq:fom-ute2}
\end{equation}

\subsubsection{MZM stability under DFD (W3i3 inheritance)}

W3i3 established that Majorana zero modes in topological superconductors are
\emph{not lifted} by DFD coupling. The DFD $\psi$-field acts as a scalar
perturbation, which preserves time-reversal symmetry for $p$-wave pairing.
The topological protection of MZMs therefore persists. For UTe$_2$:

The DFD coupling would manifest as a renormalisation of the pair potential
amplitude (through Process~A), but would not split the MZM degeneracy since
the $\psi$-field does not break the protecting symmetry (particle-hole
conjugation $\mathcal{C}$ of class DIII). The DFD-coupled UTe$_2$ remains in
the same topological class.

\begin{finding}[UTe$_2$: Topologically Stable, Thermally Impractical]
\label{finding:ute2}
UTe$_2$ achieves $\FOM_{\rm DFD} \approx 53$. Its MZM states are preserved
under DFD coupling (W3i3 result). However, the very low $\Tc = 1.6$~K makes
it impractical for P15 SRF application: operating at $T \sim 0.5$~K requires
dilution refrigeration. There is no DFD-specific advantage from the
spin-triplet pairing for Process~A: the EM stress-energy trace is insensitive
to spin structure of the order parameter.
\end{finding}

\begin{finding}[Heavy Fermions: General DFD Assessment]
\label{finding:hf-general}
All heavy fermion superconductors evaluated (CeCoIn$_5$, UBe$_{13}$, UTe$_2$)
achieve $\FOM_{\rm DFD} < 65$, far below conventional and HTS candidates.
The large effective mass enhances $\lambdaL$ (beneficial for Process~A) but
the exponential $\Tc$ penalty in surface resistance is decisive. Heavy fermion
materials are \emph{not recommended} for P15 resonator applications.
The one possible niche is a DFD-specific experiment exploiting MZM degeneracy
in UTe$_2$, but this is not relevant to the P15 generator configuration.
\end{finding}

%=============================================================================
\section{Flat-Band Materials}
\label{sec:flatband}
%=============================================================================

\subsection{Magic-angle twisted bilayer graphene (MATBG)}

MATBG inherits from W3i3 MatSci: a $2700\times$ DOS enhancement relative to
conventional graphene, arising from the van Hove singularity (VHS) at the
magic angle $\theta \approx 1.1^\circ$. The superconducting state parameters:

\begin{itemize}
\item $\Tc \approx 1$--$3$~K (carrier density dependent; maximum at VHS filling);
\item Pairing symmetry: debated ($d + id$, $p+ip$, or nematic $d$-wave);
  DFD framework predicts the $\psi$-coupling selects whichever symmetry
  maximises $\int |\psi_{\mathbf{k}}|^2 d\mathbf{k}$ on the Fermi surface
  (the symmetric channel);
\item $\lambdaL$: very large, $\sim 1$--$10~\mu$m, due to extremely low
  carrier density ($n \sim 10^{11}$~cm$^{-2}$ vs.\ $\sim 10^{22}$~cm$^{-3}$
  for metals). This is a 2D material, so $\lambdaL$ has different dimensions;
\item Pearl length $\Lambda_{\rm Pearl} = 2\lambdaL^2/d \approx 10$--$100~\mu$m
  (where $d \sim 0.7$~nm is the bilayer thickness);
\item $\xsc \approx 10$--$30$~nm.
\end{itemize}

\subsubsection{2D coupling formula for MATBG}

For a 2D superconductor, the Process~A source per unit area at fixed current
density $J_s$ is:
\begin{equation}
S_\psi^{\rm 2D} = J_s^2 \cdot \Lambda_{\rm Pearl} / (n_s^{\rm 2D} e^2),
\end{equation}
where $n_s^{\rm 2D}$ is the 2D superfluid density. The Pearl length
$\Lambda_{\rm Pearl}$ plays the role of $\lambdaL$ for 2D systems.

However, the relevant quantity for a cavity wall coating is the product
of the 2D source with the number of layers per cavity area. A single MATBG
sheet is $\sim 0.7$~nm thick. For a practical cavity coating at
$\sim 1~\mu$m total thickness (which W3i3 identified as the relevant
length scale), we would need $\sim 1000$ MATBG layers.

Stacking $N$ MATBG layers is possible but: (a) each bilayer must be
independently aligned to the magic angle with $\pm 0.05^\circ$ precision;
(b) the interlayer coupling between successive MATBG sheets would disrupt
the flat-band condition. There is no practical pathway to thick MATBG
coatings.

\subsubsection{DFD FOM for a single MATBG sheet}

For the purposes of a FOM comparison, we use the Pearl-length analogue:
\begin{equation}
\lambdaL^{\rm eff,MATBG} \equiv \Lambda_{\rm Pearl}/2 \approx 30~\mu{\rm m}
= 3\ee{4}~{\rm nm},
\end{equation}
\begin{equation}
\xsc^{\rm MATBG} \approx 20~{\rm nm},
\end{equation}
\begin{equation}
\kappa_{\rm MATBG} \approx \frac{3\ee{4}}{20} = 1500,
\end{equation}
\begin{equation}
\FOM_{\rm DFD}^{\rm MATBG}
= 1500 \times \sqrt{2}
= 1500 \times 1.41
\approx 2120.
\label{eq:fom-matbg}
\end{equation}

This is a very high FOM, but the caveat is that it applies to a single
2D sheet. The practical enhancement over a 3D material coating is suppressed
by the ratio of the sheet thickness ($0.7$~nm) to the relevant 3D length
($\lambdaL^{\rm 3D}$): the enhancement factor per unit volume is reduced by
$d/\lambdaL^{\rm 3D} \sim 10^{-2}$.

\begin{finding}[MATBG: Extraordinary FOM, Impractical at Scale]
\label{finding:matbg}
MATBG achieves an extraordinarily high $\FOM_{\rm DFD} \approx 2120$ per
sheet due to the large Pearl length, but:
(a) $\Tc \approx 1$--$3$~K requires cryogenic operation near millikelvin;
(b) scalable thick-film coatings maintaining magic-angle alignment are
not feasible;
(c) the 2D FOM is not directly comparable to 3D bulk materials without
accounting for the sheet-volume factor.
MATBG is a scientifically interesting DFD coupling system (W3i3 enhanced
DoS of $2700\times$ is real) but is not a viable P15 resonator material.
\end{finding}

\subsection{Twisted WSe$_2$ and MoS$_2$ bilayers}

Twisted transition-metal dichalcogenide (TMD) bilayers (tWSe$_2$, tMoS$_2$)
at their magic angles show correlated phases and superconductivity at
$\Tc \lesssim 3$~K (tWSe$_2$: $\sim 3$~K, tMoS$_2$: $\sim 1$~K).

The flat-band enhancement is similar to MATBG but typically smaller DOS
amplification ($\sim 500\times$ vs.\ $2700\times$ in MATBG). The Pearl
length for tWSe$_2$ is estimated at $\Lambda_{\rm Pearl} \approx 5$--$15~\mu$m.

Effective FOM:
\begin{equation}
\kappa_{\rm tWSe_2} \approx \frac{5000~\text{nm}}{15~\text{nm}} = 333,
\end{equation}
\begin{equation}
\FOM_{\rm DFD}^{\rm tWSe_2}
\approx 333 \times \sqrt{2.5}
\approx 527.
\end{equation}

Similar practical barriers as MATBG apply. TMD bilayers are not viable for
P15 cavity applications but confirm that flat-band physics is DFD-relevant.

%=============================================================================
\section{Full Material FOM Comparison Table}
\label{sec:table}
%=============================================================================

Table~\ref{tab:fom} presents the complete DFD FOM comparison across all
evaluated materials, including W3i3 candidates (Nb, Nb$_3$Sn, YBCO,
La$_3$Ni$_2$O$_7$) and W3i4 new materials. All FOMs use
$\FOM_{\rm DFD} = \kappa \cdot \Tc^{1/2}$ (Eq.~\ref{eq:fom-dfd}), with
$\kappa = \lambdaL / \xsc$.

\begin{table}[H]
\centering
\caption{DFD Process~A Figure of Merit: all evaluated materials.
  $\FOM_{\rm DFD} = (\lambdaL/\xsc)\cdot\Tc^{1/2}$.
  ``Practical'' column: Y = usable for P15 cavity walls; N = not viable.
  For 2D materials, $\lambdaL^{\rm eff} = \Lambda_{\rm Pearl}/2$ and
  the FOM is marked with $\dagger$ (not directly comparable to 3D).}
\label{tab:fom}
\small
\begin{tabular}{lcccccccc}
\toprule
Material & $\Tc$ (K) & $\lambdaL$ (nm) & $\xsc$ (nm)
  & $\kappa$ & Pairing & $\FOM_{\rm DFD}$ & Rank & Practical \\
\midrule
\multicolumn{9}{l}{\textit{W3i3 candidates (carried forward)}} \\
Nb            & 9.2  & 40   & 38   & 1.05 & $s$      & 3.2    & 10 & Y \\
Nb$_3$Sn      & 18   & 65   & 3    & 21.7 & $s$      & 92     & 5  & Y \\
YBCO          & 93   & 150  & 1.5  & 100  & $d_{x^2-y^2}$ & 965 & 2 & Y \\
La$_3$Ni$_2$O$_7$ & 80 & 100 & 3.15 & 31.7 & $d+s_\pm$ & 284 & 4 & N$^a$ \\
\midrule
\multicolumn{9}{l}{\textit{W3i4 new materials}} \\
MgB$_2$       & 39   & 85   & 4.25 & 20.0 & $s$ (two-band) & 125 & 6 & Y \\
NdNiO$_2$     & 12   & 230  & 4    & 57.5 & $d_{x^2-y^2}$ & 199 & 3 (practical: N$^b$) & N \\
CeCoIn$_5$    & 2.3  & 240  & 6    & 40   & $d_{x^2-y^2}$ & 61 & 8  & N$^c$ \\
UBe$_{13}$    & 0.86 & 800  & 15   & 53   & unknown  & 49     & 9  & N$^c$ \\
UTe$_2$       & 1.6  & 500  & 12   & 42   & $p$-wave & 53     & 7  & N$^c$ \\
MATBG$^\dagger$ & 2  & 30000 & 20  & 1500 & $d+id$?  & 2120$^\dagger$ & 1$^\dagger$ & N$^d$ \\
tWSe$_2$$^\dagger$ & 2.5 & 5000 & 15 & 333 & unknown & 527$^\dagger$ & ---$^\dagger$ & N$^d$ \\
\midrule
\multicolumn{9}{l}{\textit{Practical ranking (viable materials only)}} \\
YBCO          & 93   & 150  & 1.5  & 100  & $d_{x^2-y^2}$ & \textbf{965} & \textbf{P1} & Y \\
MgB$_2$       & 39   & 85   & 4.25 & 20.0 & $s$      & \textbf{125} & \textbf{P2} & Y \\
Nb$_3$Sn      & 18   & 65   & 3    & 21.7 & $s$      & \textbf{92}  & \textbf{P3} & Y \\
Nb            & 9.2  & 40   & 38   & 1.05 & $s$      & \textbf{3.2} & \textbf{P4} & Y \\
\bottomrule
\end{tabular}
\vspace{1mm}

\noindent\small
$^a$ Requires $> 14$~GPa pressure; no ambient-pressure practical route for cavity fabrication.\\
$^b$ Only available as thin films $< 20$~nm on specific substrates; no bulk or thick-film synthesis.\\
$^c$ $\Tc < 2.5$~K: requires dilution refrigeration; impractical for P15 sustained operation.\\
$^d$ 2D material; FOM not directly comparable; no practical thick-coating route.
\end{table}

%=============================================================================
\section{Process A Coupling: Detailed Quantitative Estimates}
\label{sec:quant}
%=============================================================================

\subsection{Absolute $\psi$-sourcing amplitude}

For a TESLA-style SRF cavity ($V_{\rm cav} = 0.01$~m$^3$, $\omega/2\pi = 1.3$~GHz,
$B_{\rm peak} = 200$~mT, $Q_0 = 10^{10}$), the Process~A $\psi$-amplitude
at the cavity centre was established in W3i3 as:
\begin{equation}
\psi_0^{\rm Nb} = C_\psi^{\rm Nb} \cdot B_{\rm peak}^2 \cdot \lambdaL^{\rm Nb}
\approx 1.44\ee{-10} \cdot 0.04 \cdot 40 \approx 2.3\ee{-10}~[\text{DFD units}].
\end{equation}

For other materials at the same $B_{\rm peak}$, the scaling is purely through
$\lambdaL$:
\begin{align}
\psi_0^{\rm MgB_2} &= \psi_0^{\rm Nb} \times \frac{\lambdaL^{\rm MgB_2}}{\lambdaL^{\rm Nb}}
  = 2.3\ee{-10} \times \frac{85}{40} = 4.9\ee{-10}, \\
\psi_0^{\rm Nb_3Sn} &= 2.3\ee{-10} \times \frac{65}{40} = 3.7\ee{-10}, \\
\psi_0^{\rm YBCO} &= 2.3\ee{-10} \times \frac{150}{40} = 8.6\ee{-10}.
\end{align}

The Nb$_3$Sn + YBCO laminate (W3i3 recommendation) benefits from the YBCO
surface layer having $3.75\times$ the Nb Process~A efficiency, while the
Nb$_3$Sn substrate provides structural support and limits field penetration.

\subsection{Could MgB$_2$ replace Nb$_3$Sn in the laminate?}

The W3i3 recommendation was:
\begin{equation}
\text{Optimal resonator: Nb}_3\text{Sn substrate (bulk) + YBCO surface coating.}
\end{equation}

The role of Nb$_3$Sn in the laminate is to:
\begin{enumerate}
\item Provide a high-$B_{c2}$ substrate ($B_{c2}^{\rm Nb_3Sn} \approx 24$~T)
  that can withstand the high peak fields;
\item Deliver low RF surface resistance at 4~K;
\item Act as a base layer onto which the YBCO coating is deposited.
\end{enumerate}

MgB$_2$ has $B_{c2}^{ab} \approx 16$--$18$~T, lower than Nb$_3$Sn, and
$R_s^{\rm MgB_2}$ at 4~K is comparable to Nb$_3$Sn. Substituting MgB$_2$ for
Nb$_3$Sn would:
\begin{itemize}
\item Reduce $B_{c2}$ tolerance by $\sim 30$\%, requiring lower peak fields
  (offset by the higher $\Tc$ allowing operation at 15--20~K);
\item Increase the base-layer $\lambdaL$ from 65~nm to 85~nm, marginally
  improving Process~A from the substrate layer;
\item Enable closed-cycle cooling (15--20~K) vs.\ liquid helium (4~K) for
  the Nb$_3$Sn substrate.
\end{itemize}

The net trade-off favours MgB$_2$ only in scenarios where cryogen-free
operation is a hard requirement, and where peak field $< 16$~T is acceptable.

%=============================================================================
\section{Updated Optimal Resonator Recommendation}
\label{sec:recommendation}
%=============================================================================

\subsection{Ranking summary for practical materials}

Of the materials evaluated across W3i1--W3i4, the practically viable candidates
rank as follows (combining FOM with fabricability):

\begin{center}
\begin{tabular}{clccc}
\toprule
Rank & Material & $\FOM_{\rm DFD}$ & Fabricability & P15 viability \\
\midrule
1 & YBCO thin film (surface) & 965 & Moderate (PLD/sputtering) & High \\
2 & MgB$_2$ (bulk/film) & 125 & Good (HPCVD) & High (cryo-free) \\
3 & Nb$_3$Sn (substrate) & 92 & Good (vapour diffusion) & High \\
4 & Nb (baseline) & 3.2 & Excellent (standard SRF) & Proven \\
\bottomrule
\end{tabular}
\end{center}

\subsection{W3i4 recommendation}

\begin{mdframed}[backgroundcolor=yellow!10]
\textbf{W3i4 Updated Resonator Recommendation}

The W3i3 recommendation of \textbf{Nb$_3$Sn substrate + YBCO surface
coating} remains optimal for maximum DFD Process~A coupling, confirmed by
the full FOM analysis across all W3i4 alternative materials.

No new material from W3i4 displaces this recommendation:
\begin{itemize}
\item La$_3$Ni$_2$O$_7$ (W3i3, FOM~284) requires high pressure: impractical;
\item NdNiO$_2$ (FOM~199) is thin-film only: impractical;
\item MATBG (FOM~2120 per sheet) is 2D: impractical for cavity walls;
\item Heavy fermions: FOM~49--61, impractical temperatures.
\end{itemize}

\textbf{Amendment}: A \textbf{MgB$_2$ variant} is recommended as the
preferred alternative for deployments where:
\begin{enumerate}
\item Cryogen-free (closed-cycle) operation at 15--20~K is required;
\item Peak fields below 16~T are acceptable;
\item Manufacturing simplicity (HPCVD at $\sim 650$~$^\circ$C) is prioritised
  over maximum Process~A coupling.
\end{enumerate}

In such cases, the optimal configuration is:
\begin{equation}
\boxed{\text{MgB}_2\text{ thin film (1--3~$\mu$m) on Cu cavity substrate,}
\text{ + YBCO surface coating (100--200~nm)}}
\end{equation}
This achieves $\FOM_{\rm DFD}$ intermediate between Nb$_3$Sn+YBCO and
pure YBCO, with closed-cycle cooling compatibility.
\end{mdframed}

\subsection{Justification for W3i3 recommendation stability}

The analysis confirms that the W3i3 Nb$_3$Sn + YBCO laminate recommendation
is robust. The reasons no W3i4 material displaces it:

\begin{enumerate}
\item \textbf{YBCO remains highest practical FOM}: $\FOM_{\rm DFD} = 965$, with
  ambient-pressure synthesis and established thin-film deposition on SRF
  substrates (PLD and sputtering at CERN, Cornell, and Jefferson Lab scales);

\item \textbf{No phonon-mediated SC exceeds MgB$_2$ at $\Tc < 40$~K}: MgB$_2$
  is the best conventional superconductor below HTS temperatures, confirmed;

\item \textbf{Cuprates (YBCO, BiSSCO) remain the unique practical HTS option}:
  ambient-pressure, established fabrication, $\Tc \approx 93$--$110$~K gives
  exponentially low $R_s$;

\item \textbf{Unconventional materials all have practical barriers}: pressure,
  thin-film-only, or millikelvin cooling.
\end{enumerate}

%=============================================================================
\section{Cross-References to P2, P11, and MatSci Module}
\label{sec:xref}
%=============================================================================

\subsection{Cross-reference to P2 (Resonators)}

P2 uses resonators for $\psi$-mode detection via parametric amplification.
The W3i3 finding that psi-to-photon transduction efficiency is $\sim 10^{-35}$
means the resonator material choice affects P2 \emph{only through the
cavity $Q$ and stored energy}. For P2:

\begin{itemize}
\item YBCO-coated resonators maintain $Q_0 \sim 10^{10}$ at $T \sim 77$~K
  (liquid nitrogen), superior to Nb at the same cooling cost;
\item MgB$_2$-coated resonators at 15--20~K achieve $Q_0 \sim 10^9$--$10^{10}$,
  viable for P2 if cryo-free operation is required;
\item For P2 squeezed-mode generation: the parametric amplifier gain scales
  with $Q_0$, so YBCO remains preferred.
\end{itemize}

\subsection{Cross-reference to P11 (SRF Cavities)}

P11 covers the SRF cavity design per se. The W3i4 analysis directly informs
P11 material selection:

\begin{enumerate}
\item \textbf{Primary recommendation} (max DFD coupling): Nb$_3$Sn substrate
  + YBCO coating (W3i3, confirmed W3i4);
\item \textbf{Alternative} (cryo-free): MgB$_2$ on Cu + YBCO (W3i4 new);
\item \textbf{Proven baseline}: Nb (TESLA standard), FOM~3.2 but unparalleled
  reliability;
\item \textbf{Reject for P11}: La$_3$Ni$_2$O$_7$ (pressure), NdNiO$_2$
  (thin-film only), all heavy fermions (millikelvin), all flat-band 2D
  materials (not scalable).
\end{enumerate}

\subsection{Cross-reference to MatSci module (W3i3)}

The W3i3 MatSci module established:
\begin{itemize}
\item La$_3$Ni$_2$O$_7$ FOM$_{\rm RT} = 252$~K (using the W3i3 MatSci FOM
  definition $\Tc \cdot \lambdaL^{-2} \cdot \xsc^{-1}$). In the W3i4 DFD
  Process~A FOM ($\kappa \cdot \Tc^{1/2}$), La$_3$Ni$_2$O$_7$ scores 284,
  consistent with its position between Nb$_3$Sn and YBCO in both rankings;
\item MATBG 2700$\times$ DOS enhancement: confirmed DFD-specific, but not
  translatable to bulk cavity performance;
\item Nb$_3$Sn + YBCO laminate: confirmed optimal in both W3i3 MatSci and
  W3i4 P15 analyses;
\item Topological SCs (MZM not lifted by DFD): inherited in W3i4 UTe$_2$
  analysis; conclusion unchanged.
\end{itemize}

%=============================================================================
\section{Summary and Conclusions}
\label{sec:summary}
%=============================================================================

\subsection{FOM formula}

The DFD Process~A Figure of Merit for SRF cavity materials is:
\begin{equation}
\FOM_{\rm DFD} = \kappa \cdot \Tc^{1/2} = \frac{\lambdaL}{\xsc} \cdot \Tc^{1/2}.
\end{equation}

\subsection{Key findings}

\begin{finding}[W3i4 Material Ranking for P15]
\label{finding:ranking}
The practical material ranking for P15 DFD Process~A coupling (all materials
that can be fabricated as cavity walls) is:
\begin{enumerate}
\item \textbf{YBCO}: $\FOM_{\rm DFD} = 965$. Optimal surface coating.
\item \textbf{MgB$_2$}: $\FOM_{\rm DFD} = 125$. Best cryo-free alternative.
\item \textbf{Nb$_3$Sn}: $\FOM_{\rm DFD} = 92$. Best substrate material.
\item \textbf{Nb}: $\FOM_{\rm DFD} = 3.2$. Proven baseline.
\end{enumerate}
All W3i4 new materials (NdNiO$_2$, heavy fermions, flat-band 2D systems)
have either higher FOM but are impractical, or have both lower FOM and
are impractical. No W3i4 material enters the practical top-4 ranking.
\end{finding}

\begin{finding}[Resonator Recommendation Confirmed]
\label{finding:resonator}
The W3i3 recommendation of \textbf{Nb$_3$Sn + YBCO laminate} as the
optimal P15/P11 resonator material is \emph{confirmed} by the W3i4 analysis.
A new secondary recommendation is added: \textbf{MgB$_2$ + YBCO} for
cryo-free ($15$--$20$~K) operation.
\end{finding}

\begin{finding}[Heavy Fermions: Not Viable for P15]
\label{finding:hf-final}
All heavy fermion superconductors (CeCoIn$_5$, UBe$_{13}$, UTe$_2$) achieve
$\FOM_{\rm DFD} < 65$ and require millikelvin cooling. They are not viable
for P15 SRF generator applications.
\end{finding}

\begin{finding}[Flat-Band 2D Materials: DFD-Relevant but Not Scalable]
\label{finding:flatband-final}
MATBG and twisted TMDs show extraordinary per-sheet FOM values ($> 500$)
due to large Pearl lengths, but cannot be scaled to cavity-wall dimensions.
Their DFD-specific DOS enhancement is real (W3i3: $2700\times$ for MATBG)
but does not translate to practical P15 enhancement. A dedicated experiment
using a 2D superconductor as a DFD $\psi$-coupling probe (not a cavity
wall) remains scientifically interesting.
\end{finding}

\begin{finding}[W3i4 Process A Coupling Quantification]
\label{finding:processa}
At $B_{\rm peak} = 200$~mT, the Process~A $\psi$-sourcing amplitude for
the Nb$_3$Sn + YBCO laminate is $8.6\ee{-10}$ (DFD units) at the YBCO
surface, compared to $2.3\ee{-10}$ for pure Nb. This $3.75\times$ enhancement
is the surface-mechanism advantage identified in W3i2 and confirmed in W3i3.
No W3i4 material offers a larger \emph{practical} Process~A enhancement.
\end{finding}

\subsection{Open questions for W3i5}

\begin{enumerate}
\item Can La$_3$Ni$_2$O$_7$ be synthesised at ambient pressure via thin-film
  routes compatible with cavity substrates? If so, its FOM~284 would enter
  the practical ranking between YBCO and Nb$_3$Sn.
\item What is the optimal YBCO coating thickness on the MgB$_2$/Cu laminate
  for maximum Process~A at 15--20~K?
\item Can the MATBG DOS enhancement ($2700\times$) be accessed via a
  surface-mounted probe geometry in P15, rather than a cavity-wall geometry?
\item Experimental: what does MgB$_2$ RF surface resistance at
  $\omega/2\pi = 1.3$~GHz and $T = 15$~K look like? Literature data at
  this frequency is sparse.
\end{enumerate}

%=============================================================================
\section*{Appendix: Numerical Cross-Check of FOM Values}
\label{app:check}
%=============================================================================

All FOM calculations are verified numerically below for transparency.

\begin{center}
\begin{tabular}{lccccc}
\toprule
Material & $\lambdaL$ (nm) & $\xsc$ (nm) & $\kappa$ & $\Tc^{1/2}$ (K$^{1/2}$)
  & $\FOM_{\rm DFD}$ \\
\midrule
Nb            & 40   & 38   & 1.053 & 3.033 & 3.19   \\
Nb$_3$Sn      & 65   & 3.0  & 21.67 & 4.243 & 91.9   \\
YBCO          & 150  & 1.5  & 100.0 & 9.644 & 964.4  \\
La$_3$Ni$_2$O$_7$ & 100 & 3.15 & 31.75 & 8.944 & 283.9 \\
MgB$_2$       & 85   & 4.25 & 20.00 & 6.245 & 124.9  \\
NdNiO$_2$     & 230  & 4.0  & 57.50 & 3.464 & 199.2  \\
CeCoIn$_5$    & 240  & 6.0  & 40.00 & 1.517 & 60.7   \\
UBe$_{13}$    & 800  & 15.0 & 53.33 & 0.927 & 49.4   \\
UTe$_2$       & 500  & 12.0 & 41.67 & 1.265 & 52.7   \\
MATBG$^\dagger$ & 30000 & 20.0 & 1500 & 1.414 & 2121$^\dagger$ \\
\bottomrule
\end{tabular}
\end{center}

\noindent$^\dagger$ 2D material; $\lambdaL^{\rm eff} = \Lambda_{\rm Pearl}/2$;
not comparable to 3D values.

\end{document}
